% ===== PRÉAMBULE CANONIQUE — à coller VERBATIM en tête de CHAQUE .tex =====
\documentclass[11pt,a4paper]{article}
\usepackage[utf8]{inputenc}\usepackage[T1]{fontenc}\usepackage{lmodern}
\usepackage[french]{babel}
\usepackage{amsmath,amssymb,mathtools}\usepackage{icomma}
\usepackage{siunitx}\sisetup{locale=FR}  % \qty{}{}, \si{}, \ang{}
\usepackage{xcolor}\usepackage{colortbl}
\usepackage{tikz}\usepackage{pgfplots}\pgfplotsset{compat=1.18}
\usepackage[siunitx,europeanresistors]{circuitikz} % circuits (élec.)
\usepackage[most]{tcolorbox}\usepackage{enumitem}\usepackage{array,booktabs}
\usepackage[a4paper,margin=2.2cm]{geometry}
\usetikzlibrary{arrows.meta,calc,angles,quotes,positioning,patterns,%
  backgrounds,shapes.geometric,decorations.pathmorphing,decorations.markings,intersections}
\definecolor{primary}{RGB}{222,49,99}\definecolor{primarydark}{RGB}{150,22,55}
\definecolor{defcol}{RGB}{0,114,178}\definecolor{methcol}{RGB}{52,92,148}
\definecolor{excol}{RGB}{0,135,90}\definecolor{attncol}{RGB}{200,40,55}
\tcbset{boxbase/.style={enhanced,breakable,boxrule=0.9pt,arc=2.5pt,
  left=3mm,right=3mm,top=2.3mm,bottom=2.3mm,fonttitle=\bfseries,coltitle=white}}
% --- boîtes TUTORIEL (noms EXACTS, ne pas renommer) ---
\newtcolorbox{cle}[1][]{boxbase,colback=defcol!6,colframe=defcol,
  title={\textsc{À retenir}\ifx\relax#1\relax\relax\else\textmd{ : \itshape #1}\fi}}
\newtcolorbox{methode}[1][]{boxbase,colback=methcol!7,colframe=methcol,sharp corners,
  title={\textsc{Méthode}\ifx\relax#1\relax\relax\else\textmd{ --- \itshape #1}\fi}}
\newtcolorbox{exemplebox}[1][]{boxbase,colback=excol!5,colframe=excol,
  title={\textsc{Exemple}\ifx\relax#1\relax\relax\else\textmd{ : \itshape #1}\fi}}
\newtcolorbox{piege}[1][]{boxbase,colback=attncol!6,colframe=attncol,
  title={\textsc{Piège}\ifx\relax#1\relax\relax\else\textmd{ : \itshape #1}\fi}}
\newtcolorbox{remarque}[1][]{boxbase,colback=black!4,colframe=black!45,
  title={\textsc{Remarque}\ifx\relax#1\relax\relax\else\textmd{ : \itshape #1}\fi}}
% --- boîtes EXERCICE / CORRIGÉ (noms EXACTS, auto-comptées) ---
\newtcolorbox[auto counter]{exo}[1][]{boxbase,colback=primary!4,colframe=primary,
  title={\textsc{Exercice}~\thetcbcounter\ifx\relax#1\relax\relax\else\textmd{ --- \itshape #1}\fi}}
\newtcolorbox[auto counter]{corrige}[1][]{boxbase,colback=excol!4,colframe=excol,
  title={\textsc{Corrigé}~\thetcbcounter\ifx\relax#1\relax\relax\else\textmd{ --- \itshape #1}\fi}}
\newcommand{\vv}[1]{\vec{#1}}\newcommand{\ds}{\displaystyle}
\newcommand{\R}{\mathbb{R}}\newcommand{\N}{\mathbb{N}}\newcommand{\Z}{\mathbb{Z}}
% =========================================================================
\begin{document}
% THEME_TITLE: Associations de résistances et générateur réel
\usetikzlibrary{babel}

\begin{center}
{\LARGE\bfseries\color{primarydark} Thème 5 — Associations de résistances et générateur réel}\\[2pt]
{\color{primary}\itshape Série, parallèle, résistance équivalente, et modèle $U=E-rI$}
\end{center}

\begin{cle}[association en série]
Des résistances \textbf{en série} sont parcourues par la \emph{même} intensité ; leurs tensions
s'ajoutent. La résistance équivalente est la \textbf{somme} :
\[ \boxed{\;R_{\text{eq}}=R_1+R_2+\cdots+R_n\;}\qquad (R_{\text{eq}}>\text{chacune}). \]
\end{cle}

\begin{center}
\begin{circuitikz}[scale=0.85,transform shape,>=Stealth]
  \draw (0,0) to[short,o-] (0.5,0) to[R,l=$R_1$] (2.3,0) to[R,l=$R_2$] (4.1,0)
        to[R,l=$R_3$] (5.9,0) to[short,-o] (6.4,0);
  \draw[defcol,<->,line width=0.9pt] (0.5,0.85) -- (5.9,0.85) node[midway,above,font=\footnotesize,defcol]{$U_{AB}$};
  \draw[->,primary,line width=0.9pt] (0.7,-0.5) -- (1.3,-0.5) node[midway,below,font=\footnotesize,primary]{$I$};
\end{circuitikz}
\end{center}

\begin{cle}[association en parallèle]
Des résistances \textbf{en parallèle} ont la \emph{même} tension ; leurs intensités s'ajoutent (loi
des nœuds). Les \emph{inverses} s'ajoutent :
\[ \boxed{\;\dfrac{1}{R_{\text{eq}}}=\dfrac{1}{R_1}+\cdots+\dfrac{1}{R_n}\;}\qquad (R_{\text{eq}}<\text{la plus petite}). \]
La \textbf{conductance} $G=1/R$ (en \si{\siemens}) s'ajoute : $G_{\text{eq}}=G_1+\cdots+G_n$.
\end{cle}

\begin{center}
\begin{circuitikz}[scale=0.8,transform shape,>=Stealth]
  \coordinate (A) at (0,0); \coordinate (B) at (4.4,0);
  \draw[thick] (A) -- (0,1.4); \draw (0,1.4) to[R,l=$R_1$] (4.4,1.4); \draw[thick] (4.4,1.4)--(B);
  \draw (A) to[R,l=$R_2$] (B);
  \draw[thick] (A) -- (0,-1.4); \draw (0,-1.4) to[R,l=$R_3$] (4.4,-1.4); \draw[thick] (4.4,-1.4)--(B);
  \draw[->,primary,line width=1pt] (-1.4,0) -- (-0.12,0) node[midway,above,font=\footnotesize,primary]{$I$};
  \fill (A) circle(1.6pt); \node[below left,font=\footnotesize] at (A){$A$};
  \fill (B) circle(1.6pt); \node[below right,font=\footnotesize] at (B){$B$};
\end{circuitikz}
\end{center}

\begin{cle}[deux astuces qui font gagner du temps]
\begin{itemize}[leftmargin=1.5em,topsep=1pt,itemsep=0pt]
  \item \textbf{Deux} résistances en parallèle : $R_{\text{eq}}=\dfrac{R_1\,R_2}{R_1+R_2}$ (produit sur somme).
  \item \textbf{$n$ résistances égales} $R$ en parallèle : $R_{\text{eq}}=\dfrac{R}{n}$.
\end{itemize}
\end{cle}

\begin{cle}[générateur réel et moteur]
Tout générateur possède une \textbf{résistance interne} $r$ : la tension à ses bornes chute quand il
débite : \[ \boxed{\;U=E-r\,I\;}. \]
$E$ = f.é.m. = tension \textbf{à vide} ($I=0$). En \textbf{court-circuit} ($U\approx0$) :
$I_{\text{cc}}=\dfrac{E}{r}$. Pour un \textbf{moteur} : $U=E'+r\,I$ (il faut fournir plus que la f.c.é.m. $E'$).
\end{cle}

\begin{center}
\begin{tikzpicture}
\begin{axis}[width=8.5cm,height=3.6cm,axis lines=middle,xlabel={$I$ (A)},ylabel={$U$ (V)},
  xmin=0,xmax=5.2,ymin=0,ymax=7,xtick={0,1,2,3,4,5},ytick={0,2,4,6},
  tick label style={font=\scriptsize},label style={font=\footnotesize}]
  \addplot[defcol,very thick,domain=0:5.5,samples=2]{6-x};
  \addplot[only marks,mark=*,defcol] coordinates {(0,6)};
  \node[defcol,font=\scriptsize,anchor=west] at (axis cs:1.6,5.5){$E$ = tension à vide ($I=0$)};
  \node[defcol,font=\scriptsize] at (axis cs:3.7,3.1){$U=E-rI$ (pente $-r$)};
\end{axis}
\end{tikzpicture}
\end{center}

\begin{methode}[simplifier un réseau de résistances]
\begin{enumerate}[leftmargin=1.7em,topsep=1pt,itemsep=0pt]
  \item Repérer les \emph{blocs} purement série ou purement parallèle et les remplacer par leur $R_{\text{eq}}$.
  \item Recommencer sur le circuit simplifié jusqu'à une seule résistance.
  \item Appliquer $I=U/R_{\text{eq}}$, puis remonter aux tensions/intensités de chaque branche.
\end{enumerate}
\end{methode}

\begin{piege}[le piège du parallèle]
En parallèle, c'est $\dfrac{1}{R_{\text{eq}}}$ qui est la somme des $\dfrac{1}{R_i}$ — \textbf{ne pas
oublier d'inverser} à la fin. La $R_{\text{eq}}$ parallèle est toujours \emph{plus petite} que la plus
petite des résistances.
\end{piege}

\end{document}
